\subsection{User Study}\label{subsec:userstudy}

Public XR SDKs differ substantially in their interfaces, including coordinate-frame conventions, joint layouts, and tightly coupled retargeting logic. These differences often require component-specific adjustments and fragmented setup procedures, increasing integration burden during avatar authoring.
To assess whether a unified authoring pattern mitigates these challenges, we conducted a controlled user study comparing MotionDAG with a public SDK baseline. We hypothesized that MotionDAG’s consistent authoring abstraction would reduce in-editor configuration time and improve perceived usability and workload.
The co-primary endpoints were total in-editor configuration time ($T_{\mathrm{cfg}}$), System Usability Scale (SUS)~\cite{brooke1996sus}, and unweighted NASA Task Load Index (NASA--TLX)~\cite{hart1988nasa_tlx}.

\subsubsection{Participants and Study Design}
We recruited nineteen participants ($N = 19$) with varying levels of XR and Unity development experience. On 5-point Likert scales, mean self-reported experience was 3.21 for programming and 2.84 for XR and game-engine graphics, indicating intermediate proficiency. 
All participants provided informed consent and were compensated.

The study employed a within-subjects design comparing two SDK conditions: \textbf{SDK~A} (MotionDAG) and \textbf{SDK~B} (public baseline). Participants completed equivalent authoring tasks in both conditions to reproduce predefined reference motions in a desktop XR simulator. 
The tasks included (1) \textit{Custom-Hands}, configuring hand tracking for custom meshes, and (2) \textit{Full-Body}, configuring full-body tracking and retargeting for an avatar.
Participants iterated between editor-side authoring and Unity Play-mode validation until outputs matched the provided reference clips.

In SDK~A, participants configured the pipeline by connecting interface-compatible components. 
In SDK~B, participants first configured the Meta Core SDK hand-tracking setup (including synthetic and custom hand visual configuration) for hand tracking.
They then enabled full-body tracking through \texttt{Movement SDK} by configuring body-tracking settings and completing the retargeting setup for the avatar.
To eliminate confounding effect from build and deployment latency, all trials were conducted in a controlled desktop XR simulator environment.

Sessions were performed on identical Windows~11 desktops (dual monitors, keyboard, mouse) in a lab setting.
Unity was version-locked to \texttt{6000.0.33f1}.
Two SDK configurations were installed: \textbf{SDK~A}, our MotionDAG integration layer, and \textbf{SDK~B}, the public SDK baseline (Meta Core SDK + \texttt{Movement SDK}) at version \texttt{78.0.0}. 
Both conditions were executed under the identical Meta desktop XR simulator.

\subsubsection{Procedure and Measures}\label{subsubsec:procedure-meature}
After informed consent and a brief demographics questionnaire ($\approx$10 min), participants received a scripted orientation and completed one unlogged practice trial. 
They then performed two counterbalanced conditions (SDK~A and SDK~B; AB/BA order), each comprising the \textit{Custom-Hands} and \textit{Full-Body} configurations. 
Optimization or step skipping was prohibited, and ambiguities were reported immediately. 
After each condition, SUS and NASA--TLX were administered.

To ensure equivalent task targets in the \textit{Custom-Hands} configuration, steps unrelated to interface authoring were completed prior to timing. 
In SDK~B, achieving the reference hand-retargeting quality required manual axis-weight tuning; this calibration was treated as a pre-task adjustment and excluded from $T_{\mathrm{cfg}}$ so that comparisons reflected interface-authoring effort. 
This exclusion is conservative for SDK~A because it removes a required SDK~B operation from the timed metric.

Configuration time ($T_{\mathrm{cfg}}$) was the total in-editor configuration time per SDK, starting at the first required authoring action and ending when both configurations passed execution verification; logging was automated inside Unity. 

Subjective assessments used SUS (0--100 with adjective ratings) and unweighted NASA--TLX. 
For NASA--TLX aggregation, higher values consistently indicate higher workload. 
Co-primary inference on $T_{\mathrm{cfg}}$, SUS, and NASA--TLX used two-sided paired tests with family-wise error controlled at $\alpha{=}0.05$ via Holm--Bonferroni correction. 
We report means~$\pm$~SD, 95\% CIs for paired mean differences, exact $p$ values, and Hedges’ $g_{\mathrm{av}}$~\cite{lakens2013,morris2002}. 

\begin{table}[t]
\setlength{\tabcolsep}{0.45\tabcolsep}
\renewcommand{\arraystretch}{1.75}
\scriptsize
\centering
\caption{Results of the controlled within-subject study comparing MotionDAG (SDK~A) and the public SDK baseline (SDK~B). 
Metrics include total configuration time ($T_{\mathrm{cfg}}$), System Usability Scale (SUS), and NASA--TLX workload. Values are reported as mean $\pm$ standard deviation.}
\label{tab:userstudy-results}
\begin{tabular*}{\columnwidth}{@{\extracolsep{\fill}}>{\raggedright\arraybackslash} p{29pt} rrr@{}}
\toprule
Statistic & $T_{\mathrm{cfg}}$ (s) & NASA--TLX & SUS \\
\midrule
SDK~A ($M\!\pm\!SD$) & $564.03\!\pm\!176.58$ & $30.53\!\pm\!17.55$ & $73.42\!\pm\!17.08$ \\
SDK~B ($M\!\pm\!SD$) & $972.89\!\pm\!243.36$ & $49.74\!\pm\!18.29$ & $34.47\!\pm\!13.91$ \\
95\% CI ($A{-}B$) & $[-493.25,\,-324.48]$ & $[-24.87,\,-13.55]$ & $[28.69,\,49.21]$ \\
$t(df)$ & $-10.18\,(18)$ & $-7.13\,(18)$ & $7.98\,(18)$ \\
$p$ & $<.0001$ & $<.0001$ & $<.0001$ \\
$d$ & $-1.84$ & $-1.03$ & $+2.39$ \\
\bottomrule
\multicolumn{4}{@{}p{\columnwidth}@{}}{$^{\mathrm{a}}$CI = confidence interval; $d$ = Hedges' $g_{\mathrm{av}}$ effect size. CIs, $t$, $p$, and $d$ are computed on paired differences (SDK~A $-$ SDK~B).}
\end{tabular*}
\end{table}


\subsubsection{Results}
\textbf{Table~\ref{tab:userstudy-results}} summarizes descriptive and inferential statistics.
Across all co-primary endpoints, SDK~A outperformed SDK~B, yielding lower configuration time ($T_{\mathrm{cfg}}$), lower NASA--TLX workload, and higher SUS scores.
All co-primary effects remained significant after Holm--Bonferroni correction.

As noted in \hyperref[subsubsec:procedure-meature]{Procedure and Measures}, SDK~B required an additional manual axis-weight calibration step for practical deployment. 
This step was excluded from $T_{\mathrm{cfg}}$ to preserve interface-level comparability; thus, the reported time difference represents a conservative estimate of practical setup overhead. MotionDAG therefore exceeded the 68-point ``Good'' SUS benchmark~\cite{bangor2009sus_adjective} whereas the baseline remained ``Poor.'' 

Order effects were assessed by computing within-subject differences ($\Delta{=}A{-}B$) and comparing sequence groups (AB vs.\ BA) using Welch's $t$-tests.
Sequence sensitivity was observed for configuration time ($t(16.94)=3.32$, $p=.0040$; mean $\Delta T_{\mathrm{cfg}}=-307.6$~s for AB and $-521.4$~s for BA), but not for NASA--TLX ($t(16.10)=-0.31$, $p=.763$) or SUS ($t(16.99)=1.53$, $p=.145$).

In this controlled study, MotionDAG improved authoring efficiency and perceived usability relative to the industry-representative baseline. 
The reduction in configuration time and workload suggests lower integration effort under equivalent task targets. 
While sequence effects influenced the magnitude of the time benefit, usability and workload improvements were robust across orders.
These findings support the role of a unified, interface-driven graph abstraction in reducing SDK-specific integration complexity during XR avatar authoring.